{{دیوان
| عنوان = غلاموں کی نماز
| ماخذ = ویکی ماخذ
| ماخذ_لنک = https://ur.wikisource.org/wiki/%D8%BA%D9%84%D8%A7%D9%85%D9%88%DA%BA_%DA%A9%DB%8C_%D9%86%D9%85%D8%A7%D8%B2
| مقابلہ = اقبال سائبر لائبریری (اقبال اکادمی پاکستان)
| مقابلہ_لنک = https://www.iqbalcyberlibrary.net
}}
=== (تُرکی وفدِ ہلالِ احمر لاہور میں) ===

<poem>
کہا مجاہدِ تُرکی نے مجھ سے بعدِ نماز
طویل سجدہ ہیں کیوں اس قدر تمھارے امام

وہ سادہ مردِ مجاہد، وہ مومنِ آزاد
خبر نہ تھی اُسے کیا چیز ہے نمازِ غلام

ہزار کام ہیں مردانِ حُر کو دُنیا میں
انھی کے ذوقِ عمل سے ہیں اُمتّوں کے نظام

بدن غلام کا سوزِ عمل سے ہے محروم
کہ ہے مروُر غلاموں کے روز و شب پہ حرام

طویل سجدہ اگر ہیں تو کیا تعجبّ ہے
ورائے سجدہ غریبوں کو اَور کیا ہے کام

خدا نصیب کرے ہِند کے اماموں کو
وہ سجدہ جس میں ہے مِلّت کی زندگی کا پیام!
</poem>
